% ch5.tex -- Module 6 ESSENTIALS, Chapter 5: Applications of the motor effect
% Condensed from src/mod6/ch5.tex: knowledge points only.
\chapter[Motors, Generators and Braking]{Applications of the Motor Effect: Motors, Generators and Braking}\label{ch5:chap}
\covers{M6.IQ4.01, M6.IQ4.02, M6.IQ4.03}

% ---- chapter-local TikZ helpers (prefix chfive, as in the full book) ----
\tikzset{
  chfiveIn/.pic={\draw[hscink, line width=0.7pt, fill=white] (0,0) circle (0.16);
    \draw[hscink, line width=0.7pt] (-0.11,-0.11) -- (0.11,0.11) (-0.11,0.11) -- (0.11,-0.11);},
  chfiveOut/.pic={\draw[hscink, line width=0.7pt, fill=white] (0,0) circle (0.16);
    \fill[hscink] (0,0) circle (0.05);},
  chfiveWire/.style={hscink, line width=1.4pt},
  pics/chfivePoles/.style={code={
    \fill[hscforce!15, draw=hscink, line width=0.5pt] (-2.05,-1.15) rectangle (-1.8,1.15);
    \node[font=\small\bfseries] at (-1.925,1.38) {N};
    \fill[hscvel!15, draw=hscink, line width=0.5pt] (1.8,-1.15) rectangle (2.05,1.15);
    \node[font=\small\bfseries] at (1.925,1.38) {S};
    \foreach \yy in {-0.95,-0.35,0.35,0.95}
      \draw[field, hscfield!45] (-1.8,\yy) -- (1.8,\yy);}},
}

\begin{summarybox}{In this chapter}
\begin{itemize}
  \item The torque on a coil, and the job of each part of a DC motor.
  \item Back emf: the large start-up current and the top speed of a motor.
  \item AC and DC generators, AC induction motors and magnetic brakes.
  \item Lenz's law as conservation of energy in all of these devices.
\end{itemize}
\end{summarybox}

%% ------------------------------------------------------------------
\section{Torque on a coil and the simple DC motor}\label{ch5:sec:torque}
\covers{M6.IQ4.01}

\begin{bigidea}
A current-carrying coil in a magnetic field feels opposite forces on its two sides, so
it turns. The torque is largest when its plane is parallel to the field.
\end{bigidea}

\begin{keywords}
  \kw{Normal}{a line at \ang{90} to the flat face (the plane) of the coil.}
  \kw{Torque ($\tau$)}{the turning effect of a force. Unit: newton metre (\si{N.m}).}
  \kw{Stator}{the part of a motor that does not move. It makes the magnetic field.}
  \kw{Rotor (armature)}{the part that turns: a coil on a soft-iron core, fixed to the axle.}
  \kw{Split-ring commutator}{two half-rings that reverse the current every half turn.}
  \kw{Brushes}{fixed carbon blocks that connect the spinning coil to the supply.}
\end{keywords}

\begin{itemize}
  \item A \term{coil} is many loops (turns) of wire. Each long side is at \ang{90} to
        $B$. With $n$ turns, it feels a force $F = nlIB$.
  \item The long sides carry opposite currents, so one force is up and one is down. They
        form a \term{couple}: equal, opposite forces, not in line, that turn an object.
  \item The forces on the short end sides point along the axle. They give no torque.
  \item The forces keep the same size. Only the \term{lever arm} $r_\perp$ (axle to the
        line of the force) changes as the coil turns (\cref{ch5:fig:endon}).
  \item Only one side in the field: the force on that side is $F = lIB$ (the side is at
        \ang{90} to $B$), and $\tau = r_\perp F$.
\end{itemize}

\begin{formulacard}{Torque on a coil}
  \cardformula{\tau = nIA_\perp B = nIAB\sin\theta}
  \begin{symbols}
    \sym{$\tau$}{torque on the coil}{\si{N.m}}
    \sym{$n$}{number of turns in the coil}{no unit}
    \sym{$I$}{current in the coil}{\si{A}}
    \sym{$A$, $A_\perp$}{area of one turn of the coil; $A_\perp = A\sin\theta$}{\si{m^2}}
    \sym{$B$}{magnetic field strength}{\si{T}}
    \sym{$\theta$}{angle between $B$ and the normal to the coil}{degree (\si{\degree})}
  \end{symbols}
  \whentouse{a flat coil turning on an axle in a uniform field.}
  \trap{$\theta$ is from the \emph{normal}: plane at \ang{30} to $B$ means $\theta = \ang{60}$.}
\end{formulacard}

\begin{figure}
  \centering
  \begin{tikzpicture}[x=1cm, y=1cm]
    % (a) theta = 90
    \begin{scope}
      \pic {chfivePoles};
      \draw[guide, ->] (1.06,0) -- (1.62,0) node[pos=0.45, below, inner sep=1pt, font=\footnotesize, text=hscmuted] {$B$};
      \draw[->, hscmuted, line width=0.7pt, dashed] (0,0) -- (90:1.25) node[right, font=\footnotesize, text=hscmuted] {normal};
      \draw[hscmuted] (0.42,0) arc (0:90:0.42);
      \node[font=\footnotesize, text=hscmuted] at (45:0.62) {$\theta$};
      \draw[hscink, line width=2.2pt] (180:0.85) -- (0:0.85);
      \pic at (180:0.85) {chfiveOut};
      \pic at (0:0.85) {chfiveIn};
      \draw[force] ($(180:0.85)+(0,0.18)$) -- ++(0,0.62) node[left, text=hscforce, font=\footnotesize] {$F$};
      \draw[force] ($(0:0.85)+(0,-0.18)$) -- ++(0,-0.62) node[right, text=hscforce, font=\footnotesize] {$F$};
      \fill[hscink] (0,0) circle (1.6pt);
      \node[font=\footnotesize, align=center, anchor=north] at (0,-1.7)
        {(a) Plane $\parallel B$, $\theta = \ang{90}$:\\largest torque.};
    \end{scope}
    % (b) theta = 50
    \begin{scope}[xshift=4.75cm]
      \pic {chfivePoles};
      \draw[guide, ->] (0,0) -- (1.3,0) node[pos=0.85, below, inner sep=1pt, font=\footnotesize, text=hscmuted] {$B$};
      \draw[->, hscmuted, line width=0.7pt, dashed] (0,0) -- (50:1.25) node[above, font=\footnotesize, text=hscmuted] {normal};
      \draw[hscmuted] (0.42,0) arc (0:50:0.42);
      \node[font=\footnotesize, text=hscmuted] at (25:0.62) {$\theta$};
      \draw[hscink, line width=2.2pt] (140:0.85) -- (-40:0.85);
      \draw[guide] (140:0.85) -- ({0.85*cos(140)},0);
      \draw[hscacc, line width=1.4pt] ({0.85*cos(140)},0) -- (0,0);
      \node[font=\footnotesize, text=hscacc, below] at ({0.425*cos(140)},0) {$r_\perp$};
      \pic at (140:0.85) {chfiveOut};
      \pic at (-40:0.85) {chfiveIn};
      \draw[force] ($(140:0.85)+(0,0.18)$) -- ++(0,0.62) node[left, text=hscforce, font=\footnotesize] {$F$};
      \draw[force] ($(-40:0.85)+(0,-0.18)$) -- ++(0,-0.62) node[right, text=hscforce, font=\footnotesize] {$F$};
      \fill[hscink] (0,0) circle (1.6pt);
      \node[font=\footnotesize, align=center, anchor=north] at (0,-1.7)
        {(b) $\theta = \ang{50}$: smaller lever\\arm, smaller torque.};
    \end{scope}
    % (c) theta = 0
    \begin{scope}[xshift=9.5cm]
      \pic {chfivePoles};
      \draw[->, hscmuted, line width=0.7pt, dashed] (0,0) -- (0:1.25) node[pos=0.62, above=0.5pt, inner sep=0.8pt, fill=white, font=\scriptsize, text=hscmuted] {normal};
      \draw[hscink, line width=2.2pt] (90:0.85) -- (-90:0.85);
      \pic at (90:0.85) {chfiveOut};
      \pic at (-90:0.85) {chfiveIn};
      \draw[force] ($(90:0.85)+(0,0.18)$) -- ++(0,0.42) node[right, text=hscforce, font=\footnotesize] {$F$};
      \draw[force] ($(-90:0.85)+(0,-0.18)$) -- ++(0,-0.42) node[right, text=hscforce, font=\footnotesize] {$F$};
      \fill[hscink] (0,0) circle (1.6pt);
      \node[font=\footnotesize, align=center, anchor=north] at (0,-1.7)
        {(c) Plane $\perp B$ ($\theta = 0$): forces\\in line, zero torque.};
    \end{scope}
  \end{tikzpicture}
  \caption{A coil seen end-on (axle = black dot). $\odot$ = current out of the page,
    $\otimes$ = into the page. $\theta$ is between $B$ and the normal.}\label{ch5:fig:endon}
\end{figure}


\subsection{The simple DC motor}
\begin{sidebyside}[0.49]
  \centering
  \begin{tikzpicture}[x={(1cm,0cm)}, y={(0cm,1cm)}, z={(0.42cm,0.30cm)}, scale=0.72]
    \def\mz{2.8}
    % N magnet (left)
    \filldraw[draw=hscink, line width=0.5pt, fill=hscforce!28] (-2.6,-1,0) -- (-2.6,-1,\mz) -- (-2.6,1,\mz) -- (-2.6,1,0) -- cycle;
    \filldraw[draw=hscink, line width=0.5pt, fill=hscforce!14] (-3.7,1,0) -- (-2.6,1,0) -- (-2.6,1,\mz) -- (-3.7,1,\mz) -- cycle;
    \filldraw[draw=hscink, line width=0.5pt, fill=hscforce!22] (-3.7,-1,0) -- (-2.6,-1,0) -- (-2.6,1,0) -- (-3.7,1,0) -- cycle;
    \node[font=\bfseries] at (-3.15,0,0) {N};
    % field lines
    \draw[field] (-2.6,0.62,1.4) -- (2.6,0.62,1.4);
    \draw[field] (-2.6,-0.92,0.25) -- (2.6,-0.92,0.25);
    % S magnet (right)
    \filldraw[draw=hscink, line width=0.5pt, fill=hscvel!14] (2.6,1,0) -- (3.7,1,0) -- (3.7,1,\mz) -- (2.6,1,\mz) -- cycle;
    \filldraw[draw=hscink, line width=0.5pt, fill=hscvel!28] (3.7,-1,0) -- (3.7,-1,\mz) -- (3.7,1,\mz) -- (3.7,1,0) -- cycle;
    \filldraw[draw=hscink, line width=0.5pt, fill=hscvel!22] (2.6,-1,0) -- (3.7,-1,0) -- (3.7,1,0) -- (2.6,1,0) -- cycle;
    \node[font=\bfseries] at (3.15,0,0) {S};
    % axle
    \draw[hscmuted, line width=0.9pt, dashed] (0,0,-1.9) -- (0,0,3.4);
    % coil: current enters the right side (into the page), leaves by the left side
    \draw[chfiveWire, midarrow] (1.4,0,0.3) -- (1.4,0,2.3);
    \draw[chfiveWire, midarrow] (1.4,0,2.3) -- (-1.4,0,2.3);
    \draw[chfiveWire, midarrow] (-1.4,0,2.3) -- (-1.4,0,0.3);
    \draw[chfiveWire] (1.4,0,0.3) -- (0.12,0,0.3) -- (0.12,0,-0.9);
    \draw[chfiveWire] (-1.4,0,0.3) -- (-0.12,0,0.3) -- (-0.12,0,-0.9);
    \node[font=\footnotesize, text=hscink] at (1.75,-0.2,1.6) {$I$};
    \node[font=\footnotesize, text=hscink] at (-1.62,0.25,1.6) {$I$};
    % forces
    \draw[force] (1.4,-0.08,1.3) -- (1.4,-1.0,1.3) node[right, text=hscforce, font=\footnotesize] {$F$};
    \draw[force] (-1.4,0.08,1.3) -- (-1.4,1.0,1.3) node[left, text=hscforce, font=\footnotesize] {$F$};
    \node[font=\footnotesize, text=hscfield] at (2.0,0.85,1.4) {$B$};
    % commutator and brushes at z = -0.9
    \begin{scope}[shift={(0,0,-0.9)}]
      \fill[white] (0,0) circle (0.36);
      \draw[hscseriesB, line width=3pt] (100:0.27) arc (100:260:0.27);
      \draw[hscseriesA, line width=3pt] (-80:0.27) arc (-80:80:0.27);
      \filldraw[fill=hscmuted!55, draw=hscink, line width=0.4pt] (-0.78,-0.09) rectangle (-0.34,0.09);
      \filldraw[fill=hscmuted!55, draw=hscink, line width=0.4pt] (0.34,-0.09) rectangle (0.78,0.09);
      \draw[hscink, line width=0.8pt] (-0.78,0) -- (-1.5,0) -- (-1.5,-2.0) -- (-0.36,-2.0);
      \draw[hscink, line width=0.8pt] (0.78,0) -- (1.5,0) -- (1.5,-2.0) -- (0.0,-2.0);
      \draw[hscink, line width=1.6pt] (-0.36,-2.14) -- (-0.36,-1.86);
      \draw[hscink, line width=0.9pt] (0,-2.3) -- (0,-1.7);
      \node[font=\footnotesize] at (0.22,-1.65) {$+$};
      \node[font=\footnotesize] at (-0.58,-1.65) {$-$};
      \node[diagramlabel, anchor=north] at (-0.18,-2.35) {DC supply};
      \draw[hscmuted, line width=0.4pt] (-0.24,-0.22) -- (-1.75,-0.95) node[anchor=north east, diagramlabel, inner sep=1pt, align=right] {split-ring\\commutator};
      \draw[hscmuted, line width=0.4pt] (0.65,-0.09) -- (1.75,-0.95) node[anchor=north west, diagramlabel, inner sep=1pt] {brush};
    \end{scope}
    \draw[hscmuted, line width=0.4pt] (0,0,3.4) -- ++(0.6,0.5) node[anchor=south west, diagramlabel, inner sep=1pt] {axle};
    \draw[hscmuted, line width=0.4pt] (-0.4,0,2.3) -- ++(0,1.05) node[anchor=south, diagramlabel, inner sep=1pt] {coil (rotor)};
    \node[diagramlabel, anchor=south, align=center] at (-3.15,1,2.9) {magnet\\(stator)};
  \end{tikzpicture}
  \captionof{figure}{A simple DC motor. Seen from the commutator end, the coil turns
    clockwise.}\label{ch5:fig:motor}
\nextside
  \begin{itemize}
    \item A \term{DC motor} changes electrical energy into rotational kinetic energy
          (\cref{ch5:fig:motor}).
    \item The split ring reverses the current every half turn, as the plane passes
          perpendicular to $B$. So the torque always turns the coil one way.
    \item At the swap the brushes touch the gaps. \term{Rotational inertia} (the
          tendency to keep spinning) carries the coil on.
  \end{itemize}
\end{sidebyside}

\textbf{No commutator:} the coil rocks, then stops with its plane perpendicular to $B$.
\textbf{Dead spot:} a single coil stopped there has no torque, so it cannot start by
itself. Many coils at different angles give a smooth torque.

\textbf{Radial field.} Curved \term{pole pieces} (the shaped ends of the magnets) and a
soft-iron core make a
\term{radial magnetic field}: its lines point straight out from the centre, like spokes.
The field lies along the plane of the coil, so $\theta = \ang{90}$ in every position. The
torque stays at $nIAB$ for almost the whole turn. The maximum is \emph{not} bigger; the
average is.

%% ------------------------------------------------------------------
\section{Back emf, maximum speed and energy}\label{ch5:sec:backemf}
\covers{M6.IQ4.01, M6.IQ4.03}

\begin{bigidea}
A spinning motor coil is also a generator. Its induced emf, the \term{back emf}
($\varepsilon_\mathrm{back}$), opposes the supply. The \term{net emf},
$V - \varepsilon_\mathrm{back}$, pushes the current.
\end{bigidea}

\begin{itemize}
  \item A spinning coil has a changing flux, so an emf is induced (Faraday's law). By
        Lenz's law it opposes the change that caused it, so it acts against the supply.
  \item The back emf is proportional to the speed. Every motor has one, AC or DC.
\end{itemize}

\begin{sidebyside}[0.40]
  \centering
  \begin{tikzpicture}
    \begin{axis}[hsc, scale only axis, width=4.0cm, height=2.2cm,
        xmin=0, xmax=5, ymin=0, ymax=1.25, xtick=\empty, ytick={0,1},
        yticklabels={0,$\dfrac{V}{R}$}, xlabel={time from switch-on}, ylabel={current}]
      \addplot[hscseriesB, domain=0:5, samples=100] {0.1+0.9*exp(-x)};
      \node[anchor=south west, font=\footnotesize, text=hscseriesB] at (axis cs:2.0,0.2) {levels off above 0};
      \node[anchor=west, font=\footnotesize, align=left] at (axis cs:0.45,0.95)
        {largest at switch-on};
    \end{axis}
  \end{tikzpicture}
  \captionof{figure}{Current in a DC motor after switch-on. The torque has the same
    shape.}\label{ch5:fig:startup}
\nextside
  \begin{enumerate}
    \item Switch-on: the coil is still, so $\varepsilon_\mathrm{back} = 0$. $R$ is
          small, so $I = \dfrac{V}{R}$ and the torque are large.
    \item Speeding up: the back emf grows, so the current and torque fall.
    \item \term{Maximum speed}: the torque just balances friction and the load. The
          back emf is a little less than $V$, so a small current still flows
          (\cref{ch5:fig:startup}). In an ideal motor (no friction, no load), the back
          emf $\approx V$ and the current $\approx 0$.
  \end{enumerate}
\end{sidebyside}

\begin{itemize}
  \item A bigger \term{load} (whatever the motor drives) slows the motor. The back emf
        falls, so the current and torque rise.
  \item A motor \term{stalls} when it is switched on but cannot turn. With no back emf,
        $I = \dfrac{V}{R}$ is large. The heating $I^2R$ can melt the insulation (``burn
        out''). A \term{starting resistor} in series limits the start-up current.
  \item Energy: $VI = \varepsilon_\mathrm{back}I + I^2R$. Power in $=$ mechanical power
        $+$ heat in the coil. So back emf is how the motor does work. A stalled motor turns
        all of its input energy into heat.
  \item \textbf{Lenz's law is conservation of energy.} Imagine the induced emf
        \emph{helped} the supply. A faster coil would get more current and torque, so it
        would speed up without limit. Energy would be made from nothing. So the induced
        emf must oppose.
\end{itemize}

\begin{formulacard}[notonsheet]{Current in a running motor}
  \cardformula{I = \frac{V - \varepsilon_\mathrm{back}}{R}}
  \begin{symbols}
    \sym{$I$}{current in the motor coil}{\si{A}}
    \sym{$V$}{supply emf (supply voltage)}{\si{V}}
    \sym{$\varepsilon_\mathrm{back}$}{back emf induced in the spinning coil}{\si{V}}
    \sym{$R$}{resistance of the coil}{\si{\ohm}}
  \end{symbols}
  \whentouse{a running DC motor. At switch-on or when stalled,
    $\varepsilon_\mathrm{back} = 0$, so $I = \dfrac{V}{R}$.}
  \derive{$V = IR$ (on the sheet) with the net emf: $V - \varepsilon_\mathrm{back} = IR$.
    Times $I$: $VI = \varepsilon_\mathrm{back}I + I^2R$ (power in $=$ mechanical $+$ heat).}
  \trap{the back emf is zero at switch-on, not at top speed.}
\end{formulacard}

%% ------------------------------------------------------------------
\section{Generators: AC and DC}\label{ch5:sec:generators}
\covers{M6.IQ4.02, M6.IQ4.03}

\begin{bigidea}
A \term{generator} changes kinetic energy into electrical energy. It is a motor run
backwards: an outside force turns the coil, the flux changes, and an emf is induced.
\end{bigidea}

\begin{itemize}
  \item As the coil turns, $\theta$ changes, so the flux $\Phi = BA\cos\theta$ changes
        and an emf is induced (Faraday's law: $\varepsilon = -N\dfrac{\Delta\Phi}{\Delta t}$).
  \item \textbf{Plane perpendicular to $B$:} flux largest ($BA$) but not changing for an
        instant, so the emf is \emph{zero}.
  \item \textbf{Plane parallel to $B$:} flux zero but changing fastest (the sides cut
        straight across the field lines), so the emf is \emph{largest}.
  \item \textbf{AC generator:} \term{slip rings} (two full rings, each always joined to
        the same end of the coil). The coil's emf reverses every half turn, so the
        output is \term{alternating current (AC)}: it keeps reversing direction.
  \item \textbf{DC generator:} a split-ring commutator swaps the connections every half
        turn, when the emf is zero. The output is \term{direct current (DC)}: one way
        only, in pulses (\cref{ch5:fig:gengraphs}).
\end{itemize}

\begin{figure}
  \centering
  \begin{tikzpicture}
    \begin{axis}[hsc, name=fa, axis x line=middle, scale only axis, width=10.2cm, height=1.0cm,
        xmin=0, xmax=2, ymin=-1.3, ymax=1.3, xtick={0,0.25,...,2}, xticklabels={},
        ytick={-1,0,1}, yticklabels={$-BA$,0,$BA$}, ylabel={$\Phi$},
        ylabel style={rotate=-90}]
      \addplot[hscseriesA, domain=0:2, samples=161] {cos(360*x)};
    \end{axis}
    \foreach \f/\ang in {0/90, 0.125/0, 0.25/90, 0.375/0, 0.5/90} {
      \coordinate (ic) at ($(fa.north west)!\f!(fa.north east)+(0,0.95)$);
      \draw[hscink, line width=1.6pt] ($(ic)+({\ang}:0.24)$) -- ($(ic)+({\ang+180}:0.24)$);
      \fill[hscink] (ic) circle (1pt);
    }
    \node[anchor=west, font=\footnotesize, text=hscmuted] at ($(fa.north west)!0.56!(fa.north east)+(0,0.95)$)
      {$\leftarrow$ the coil seen end-on (field $B$ to the right)};
    \node[anchor=south, font=\scriptsize, text=hscmuted, align=center, text depth=0.3em] at ($(fa.north west)!0.0!(fa.north east)+(0,1.22)$) {plane\\$\perp B$};
    \node[anchor=south, font=\scriptsize, text=hscmuted, align=center, text depth=0.3em] at ($(fa.north west)!0.125!(fa.north east)+(0,1.22)$) {plane\\$\parallel B$};
    \begin{axis}[hsc, name=fb, axis x line=middle, at={(fa.south west)}, anchor=north west, yshift=-0.3cm,
        scale only axis, width=10.2cm, height=1.0cm,
        xmin=0, xmax=2, ymin=-1.3, ymax=1.3, xtick={0,0.25,...,2}, xticklabels={},
        ytick={-1,0,1}, yticklabels={$-\varepsilon_\mathrm{peak}$,0,$\varepsilon_\mathrm{peak}$},
        ylabel={emf}]
      \addplot[hscseriesB, domain=0:2, samples=161] {sin(360*x)};
      \node[anchor=west, font=\footnotesize, align=left] at (axis cs:2.05,0) {AC\\(slip rings)};
    \end{axis}
    \begin{axis}[hsc, name=fc, at={(fb.south west)}, anchor=north west, yshift=-0.3cm,
        scale only axis, width=10.2cm, height=1.0cm,
        xmin=0, xmax=2, ymin=-0.3, ymax=1.3, xtick={0,0.25,...,2},
        xticklabels={0,$T/4$,$T/2$,$3T/4$,$T$,,,,$2T$},
        ytick={0,1}, yticklabels={0,$\varepsilon_\mathrm{peak}$}, ylabel={emf}, xlabel={time}]
      \addplot[hscseriesC, domain=0:2, samples=161] {abs(sin(360*x))};
      \node[anchor=west, font=\footnotesize, align=left] at (axis cs:2.05,0.5) {DC\\(split ring)};
    \end{axis}
  \end{tikzpicture}
  \caption{Flux and output emf for a coil that starts with its plane perpendicular to $B$
    and turns at a steady rate. $T$ is the time for one turn.}\label{ch5:fig:gengraphs}
\end{figure}

\begin{itemize}
  \item \textbf{Faster turning:} bigger peak emf \emph{and} higher frequency. Twice the
        speed gives twice the peak and half the period. More turns, a stronger field or
        a bigger coil raise the peak only. The other direction flips the graph.
  \item \textbf{Energy:} with a load connected, the current in the coil feels a
        motor-effect force. By Lenz's law it opposes the rotation. The work done against
        it becomes electrical energy, so more current makes the coil harder to turn.
\end{itemize}

%% ------------------------------------------------------------------
\section{AC induction motors and magnetic braking}\label{ch5:sec:induction}
\covers{M6.IQ4.02, M6.IQ4.03}

\begin{bigidea}
When a field and metal move past each other, \term{eddy currents} (loops of current
inside the metal) are induced. By Lenz's law, their forces oppose the \emph{relative}
motion.
\end{bigidea}

\subsection{AC induction motors}
In an \term{AC induction motor} the rotor current is induced, so there are no brushes
and no commutator.
\begin{enumerate}
  \item The stator has three pairs of coils on \term{three-phase AC} (three AC supplies
        a third of a cycle apart). They peak in turn, so they make a \term{rotating
        magnetic field}: fixed strength, direction turning steadily. No coil moves.
  \item The rotor is a \term{squirrel cage}: metal bars joined at both ends by end
        rings. The moving field changes the flux through each loop of bars, so currents
        are induced (Faraday).
  \item These currents are in the field, so the bars feel forces ($F = lIB$). By Lenz's
        law the forces oppose the \emph{relative} motion: the rotor is dragged after
        the field.
\end{enumerate}
\begin{itemize}
  \item \term{Slip}: the rotor always turns a little slower than the field. If it caught
        up, there would be no flux change, no current and no force. A bigger load gives
        more slip, so more current and more torque.
  \item No brushes or commutator, so little wear and no sparking. The supply frequency
        sets the speed, so it is hard to change.
\end{itemize}

\subsection{Magnetic braking}
\begin{sidebyside}[0.45]
  \centering
  \begin{tikzpicture}[x=1.0cm, y=1.0cm]
    \filldraw[fill=hscmuted!12, draw=hscink, line width=0.9pt] (0,0) circle (1.9);
    \fill[hscink] (0,0) circle (1.5pt);
    \draw[->, hscmuted, line width=0.9pt] (125:2.15) arc (125:55:2.15);
    \node[font=\footnotesize, text=hscmuted, anchor=south] at (90:2.17) {disc turns clockwise};
    \filldraw[fill=hscfield!14, draw=hscfield, dashed, line width=0.6pt] (0.35,-0.75) rectangle (1.75,0.75);
    \foreach \xx in {0.55,1.55} \foreach \yy in {-0.42,0.42}
      \node[text=hscfield, font=\footnotesize] at (\xx,\yy) {$\times$};
    \node[font=\scriptsize, text=hscfield, anchor=east, align=right] at (0.3,0.55) {$B$ into\\the page};
    \draw[hscseriesB, line width=1.2pt, ->] (1.05,0.75) ++(-90:0.3) arc (-90:250:0.3);
    \draw[hscseriesA, line width=1.2pt, ->] (1.05,-0.75) ++(90:0.3) arc (90:-250:0.3);
    \draw[->, hscink, line width=0.9pt] (0.5,-0.15) -- (1.6,-0.15) node[pos=0.25, above, font=\scriptsize] {$I$};
    \draw[force] (1.05,-0.08) -- (1.05,0.4) node[pos=0.55, right, text=hscforce, font=\footnotesize] {$F$};
    \draw[velocity] (1.85,0.45) -- (1.85,-0.45) node[pos=0.5, right, text=hscvel, font=\footnotesize] {$v$};
    \node[font=\scriptsize, text=hscseriesB, anchor=west] at (1.42,1.28) {metal enters field};
    \node[font=\scriptsize, text=hscseriesA, anchor=west] at (1.42,-1.28) {metal leaves field};
  \end{tikzpicture}
  \captionof{figure}{A metal disc turning through a field into the page ($\times$) on its
    right side.}\label{ch5:fig:disc}
\nextside
  \term{Magnetic braking} slows moving metal with eddy currents, without touching it
  (\cref{ch5:fig:disc}).
  \begin{itemize}
    \item Metal \emph{entering} the field: flux into the page grows, so the eddy current
          is anticlockwise. Metal \emph{leaving}: clockwise.
    \item In the field both currents flow to the right. The force on them points up,
          opposite to $v$, so the disc slows.
    \item Energy: kinetic $\to$ electrical $\to$ heat in the metal.
  \end{itemize}
\end{sidebyside}

\begin{itemize}
  \item The braking force grows with speed. At zero speed there is no flux change and no
        force, so it cannot hold an object still. Friction brakes finish the stop.
  \item A stronger field or a better conductor brakes more strongly. Slots break up the
        eddy-current loops and weaken it. No contact, so no wear.
  \item Heat $Q$ on a level track $=$ kinetic energy lost:
        $Q = \tfrac{1}{2}mu^2 - \tfrac{1}{2}mv^2$ ($u$ = speed before braking, $v$ = speed
        after). If it also falls a height $\Delta h$, add $mg\Delta h$.
  \item \term{Regenerative braking}: the motor of an electric car acts as a generator
        and charges the battery. By Lenz's law the force opposes the motion, so the car
        slows, and the energy is saved.
\end{itemize}

%% ------------------------------------------------------------------
\section{Key experiments}\label{ch5:sec:experiments}
\covers{M6.IQ4.01, M6.IQ4.03}

Two practicals matter here. In the \textbf{magnetic braking practical}, a magnet falls
slowly through a copper tube. A swinging aluminium plate between magnet poles stops
quickly, but a slotted plate swings longer. Eddy currents oppose the motion, and the
energy becomes heat. The \textbf{DC motor practical} is below.

\begin{experiment}[unbreakable]{Investigating a simple DC motor}
\aim{To find the job of each part of a DC motor, and the effect of back emf.}
\method{A coil of enamelled wire (wire with a thin insulating coat) rests on paper-clip
  supports above a magnet. The coat is scraped off all round one end, but off only the
  top half of the other end. Then a lab motor's current is read while it runs, and
  again while it is held still for 1--2 seconds.}
\results{The coil spins only with one end half scraped. Lab motor: held still
  \qty{1.5}{A}; running \qty{0.30}{A}.}
\conclusion{The half-scraped end switches the current off for the half turn when the
  torque would act backwards. So, like a commutator, it keeps the coil turning one way.
  The clips act as brushes. The lab motor draws less current when it runs, because its
  back emf opposes the supply.}
\significance{It shows every part of a DC motor, and back emf, in a few minutes.}
\end{experiment}

%% ------------------------------------------------------------------
\begin{summarybox}{Remember}
\begin{itemize}
  \item $\tau = nIAB\sin\theta$, $\theta$ from the normal: largest with the plane $\parallel B$.
  \item The split ring reverses the current every half turn, so the torque always turns
        the coil one way.
  \item Back emf opposes the supply and grows with speed:
        $I = \frac{V - \varepsilon_\mathrm{back}}{R}$ (not on the sheet).
  \item Generator emf is largest with the plane $\parallel B$. Slip rings: AC; split ring: DC.
  \item Induction motor: a rotating field drags the cage rotor round, a little slower (slip).
  \item Magnetic brake: eddy currents oppose the motion; no force at zero speed.
  \item Lenz's law is conservation of energy: if induced currents helped the change,
        energy would be made from nothing.
\end{itemize}
\end{summarybox}
